\documentclass[11pt]{article}
\usepackage[margin=1in]{geometry}
\usepackage{amsmath,amssymb,amsthm}
\usepackage[T1]{fontenc}
\usepackage{lmodern}
\usepackage{microtype}
\usepackage{xcolor}
\usepackage{hyperref}
\hypersetup{colorlinks=true,linkcolor=blue!50!black,urlcolor=blue!50!black}

\newtheorem{theorem}{Theorem}
\newtheorem{lemma}{Lemma}
\newtheorem{corollary}{Corollary}
\theoremstyle{definition}
\newtheorem{remark}{Remark}

\newcommand{\Smin}{S_{64}}
\newcommand{\Pp}{P^{+}}
\newcommand{\code}[1]{\texttt{#1}}

\title{An interval for $S_{64}$, the smallest Carmichael number\\ with exactly 64 prime factors}
\author{Alexander Eastwood}
\date{September 11, 2026}

\begin{document}
\maketitle

\begin{abstract}
Let $\Smin$ denote the smallest Carmichael number with exactly $64$ distinct prime factors
(the $k=64$ entry of OEIS \href{https://oeis.org/A006931}{A006931}, currently unknown). We give
an explicit two--sided bound. For the upper bound we exhibit a $148$--digit Carmichael number
$N$ with exactly $64$ prime factors, so $\Smin \le N$. For the lower bound we prove, by a finite
exhaustion over a certified prime universe, that no $64$--factor Carmichael number is smaller than
$10^{145}$. Together,
\[
  10^{145} \;\le\; \Smin \;\le\; N \;<\; 10^{148},
\]
so $\Smin$ has $146$, $147$, or $148$ decimal digits. All arithmetic is exact; the lower bound is a
reproducible computation, verified here by three independent implementations, one of which takes its
Korselt verdicts from a frozen reference oracle that shares no code with the search.
\end{abstract}

\section{Statement}

A \emph{Carmichael number} is a composite $n$ with $b^{n}\equiv b \pmod n$ for all integers $b$.
By Korselt's criterion, $n$ is Carmichael iff it is squarefree and $(p-1)\mid(n-1)$ for every prime
$p\mid n$. Write $\omega(n)$ for the number of distinct prime factors. The sequence
$S_k=\min\{\,n : n \text{ Carmichael},\ \omega(n)=k\,\}$ is OEIS A006931; values are known for
$3\le k\le 63$, and $S_{64}$ is open. This note proves:

\begin{theorem}[Interval]\label{thm:interval}
$10^{145}\le \Smin \le N$, where
\[
  N \;=\; 8664205557705627135892897321176984054850894080429746661350141011836598584916617361518185311595941924456408326192593430956147737016566872025613568001
\]
is a $148$--digit Carmichael number with exactly $64$ prime factors. Consequently $\Smin$ has
$146$, $147$, or $148$ decimal digits.
\end{theorem}

The two bounds are proved in \S\ref{sec:upper} and \S\ref{sec:lower}. We use throughout the
following elementary consequence of Korselt's criterion.

\begin{lemma}[Pairwise admissibility]\label{lem:adm}
If $p<q$ are distinct primes dividing a Carmichael number $n$, then $p\nmid(q-1)$. Moreover $n$ is
odd.
\end{lemma}

\begin{proof}
Since $q\mid n$ we have $(q-1)\mid(n-1)$, so $n\equiv 1\pmod{d}$ for every divisor $d$ of $q-1$. If
$p\mid(q-1)$ then $n\equiv 1\pmod p$; but $p\mid n$ gives $n\equiv 0\pmod p$, forcing $p\mid 1$, a
contradiction. For oddness: if $n$ were even, an odd prime factor $p$ would make the even number
$p-1$ divide the odd number $n-1$ (as $n$ is squarefree with an odd factor, $n-1$ is even only if
$n$ is odd), which is impossible; concretely, no even number $>2$ is Carmichael.
\end{proof}

\section{Upper bound}\label{sec:upper}

The number $N$ of Theorem~\ref{thm:interval} is the product of the $64$ primes
\[
\begin{aligned}
&19,\,29,\,31,\,37,\,41,\,43,\,47,\,53,\,61,\,67,\,71,\,73,\,79,\,89,\,97,\,101,\,103,\,109,\,113,\,127,\\
&131,\,137,\,139,\,151,\,157,\,163,\,167,\,181,\,193,\,197,\,199,\,211,\,239,\,241,\,251,\,257,\,271,\,277,\,281,\,307,\\
&313,\,331,\,337,\,353,\,379,\,397,\,401,\,421,\,433,\,449,\,461,\,463,\,491,\,541,\,547,\,577,\,599,\,631,\,673,\,811,\\
&829,\,883,\,1951,\,3697.
\end{aligned}
\]
These are $64$ distinct primes; their product is $N$ (re--derive $N$ by multiplying the list rather
than trusting the printed decimal), and one checks $(p-1)\mid(N-1)$ for each. Hence $N$ is a
squarefree Carmichael number with $\omega(N)=64$, and $\Smin\le N$. Since $10^{147}\le N<10^{148}$,
the number $N$ has $148$ digits. The certificate is verified by
\code{ref/ref\_carmichael.py::verify\_certificate}, an independent, dependency--free routine
(independent Miller--Rabin primality, squarefreeness, and Korselt) sharing no code with the search
that produced $N$.

\begin{remark}
$N$ is an incumbent found by an exchange meet--in--the--middle search; it is not claimed to be
minimal. It improves a previously reported $149$--digit candidate and sits between Webster's known
neighbours $N_{63}$ ($145$ digits) and the $N_{65}$ candidate ($151$ digits), where a true $\Smin$
must lie.
\end{remark}

\section{Lower bound}\label{sec:lower}

Fix $Y=10^{145}$. We show no $64$--factor Carmichael number lies below $Y$.

\subsection{A certified prime universe}

\begin{lemma}[Cofactor lower bound]\label{lem:C}
Let $n$ be a Carmichael number with $\omega(n)=64$, and write its primes as
$p_1<p_2<\dots<p_{63}<q$. Then
\[
  D_n \;:=\; \prod_{i=1}^{63} p_i \;\ge\; C,
  \qquad
  C \;=\; 1812140595092971804464959059612576188399713525008997620073689836850955613086867158799858934921557347504370550240044086773665823716701403213549,
\]
a $142$--digit integer, independent of $q$ and of any modulus.
\end{lemma}

\begin{proof}[Proof (finite computation)]
The $63$ cofactor primes are odd (Lemma~\ref{lem:adm}) and pairwise admissible. Partition the
argument on $S$, the set of cofactor primes that are $\le 97$; there are finitely many internally
admissible such $S$ (each pair in $S$ satisfying Lemma~\ref{lem:adm}). Given $S$, the remaining
$63-|S|$ cofactor primes exceed $97$ and are each individually admissible with every element of
$S$; replacing them by the $63-|S|$ \emph{smallest} primes $>97$ that are individually admissible
with $S$ can only decrease the product, and \emph{ignoring} the mutual admissibility constraints
among those tail primes can only decrease it further. Thus for each admissible $S$ we obtain a
rigorous lower bound on $\prod_i p_i$; taking the minimum over all $407760$ admissible $S$ gives the
constant $C$. A finite prime table (odd primes $\le 10^4$) is used, and the computation checks that
enough eligible tail primes are present for every $S$, so no smaller completion is truncated away.
\end{proof}

\begin{corollary}[Universe]\label{cor:univ}
Every $64$--factor Carmichael number $n<Y$ has largest prime factor
\[
  \Pp(n)\;=\;q\;\le\;\Big\lfloor \tfrac{Y-1}{C}\Big\rfloor \;=\; 5518 .
\]
In particular every prime factor of such an $n$ lies among the $727$ odd primes $\le 5518$ (the
largest being $5507$).
\end{corollary}

\begin{proof}
$n=q\,D_n\ge q\,C$ by Lemma~\ref{lem:C}, and $n<Y$ gives $q\le(Y-1)/C$; evaluate the floor.
\end{proof}

\subsection{Exhaustion}

\begin{theorem}[Lower bound]\label{thm:lower}
There is no Carmichael number $n<10^{145}$ with $\omega(n)=64$. Hence $\Smin\ge 10^{145}$.
\end{theorem}

\begin{proof}[Proof (finite exhaustion)]
By Corollary~\ref{cor:univ} every candidate is a product of $64$ distinct primes drawn from the
universe $U$ of $727$ odd primes $\le 5518$, with product $<Y$. Enumerate all increasing
$64$--subsets of $U$ with product $<Y$, pruning a partial choice only by two necessary conditions:
(i) too few primes remain to reach $64$; and (ii) the product of the current choice times the
smallest admissible primes still available already reaches $Y$ (an exact--integer $R_{\min}$ test).
Pairwise admissibility (Lemma~\ref{lem:adm}) is enforced during descent---once $p$ is chosen, every
later prime $q$ with $p\mid(q-1)$ is discarded---which removes no genuine Carmichael candidate. At
each complete $64$--product, test Korselt.

The enumeration terminates with exactly $127092$ admissible complete products below $Y$, of which
\emph{none} is Carmichael. Therefore no $64$--factor Carmichael number is smaller than $10^{145}$.
\end{proof}

\subsection{Reproducibility}\label{sec:repro}

The lower bound was computed three independent ways, all agreeing:

\begin{enumerate}
\item A self--contained certificate (\code{s64\_lower\_bound\_certificate.py}) that recomputes $C$
and the universe, then runs \emph{two} disjoint enumerators---a binary include/exclude walk using an
incrementally maintained $\operatorname{lcm}(p_i-1)$ for the Korselt test ($327049$ nodes), and an
ordered--choice walk using direct per--prime divisibility, a separate prime generator, and no
bitmasks ($163525$ nodes). Both report $127092$ admissible products, $0$ Carmichael, and an
identical order--sensitive leaf digest
\code{sha256 = e05f62598eea784b3f60f025c6ccbf30700115ecbef42702e24154dcece32321}.
\item An independent prime--bound tool (\code{carmichael\_prime\_bound.py}) that recomputes $C$ and
$\Pp\le 5518$ from first principles.
\item A separate enumerator that takes the Korselt verdict at every one of the $127092$ leaves from
the frozen oracle \code{ref/ref\_carmichael.py} rather than from its own arithmetic, again finding
$0$ Carmichael. This enumerator was validated against the oracle's brute--force
\code{least\_with\_k} on small $k$: $k=3\!\to\!561$, $k=4\!\to\!41041$, $k=5\!\to\!825265$.
\end{enumerate}

\section{The interval}

Combining Theorem~\ref{thm:lower} (lower) with \S\ref{sec:upper} (upper) proves
Theorem~\ref{thm:interval}: $10^{145}\le\Smin\le N<10^{148}$, so $\Smin$ has $146$, $147$, or $148$
digits. This is a global statement about $\Smin$; it is \emph{not} a determination of $\Smin$, which
would require excluding every $64$--factor Carmichael number in the open interval $(10^{145},N)$.

\begin{remark}[Next step]
The same certified machinery run at $Y=10^{146}$ would either exhibit a smaller Carmichael number or
prove $\Smin\ge 10^{146}$, in the latter case narrowing the interval to $147$--$148$ digits. The
universe and the admissible--product count grow with $Y$; whether the direct enumerator suffices at
$10^{146}$ or the exchange decomposition is required is the subject of ongoing work.
\end{remark}

\section*{Acknowledgements}

The known values and neighbours for $k=3$--$63$ (and the $k=65$ candidate) are from Jonathan
Webster's \href{https://github.com/jewebste/small-carmichael-numbers}{\code{small-carmichael-numbers}}
(Butler University for $k=36$--$63$). The design of the lower--bound certificate benefited from
consultation with a large--language reasoning model (OpenAI GPT--6); the result reported here was
independently reproduced by the three implementations of \S\ref{sec:repro}, including one gated by
the frozen oracle. Existence of Carmichael numbers with every large factor count is due to
Alford--Grantham--Hayman--Shallue and Larsen--Wright; the values above are new only as bounds on the
\emph{minimum}.

\end{document}
